% =============================================================================
% Chapter 3: Data
% Plan: docs/thesis/REWRITE_PLAN.md, section 2 "3 Data". Numbers: numbers_data.tex (\Data... macros).
% Absorbs draft_2026-10_original/main.tex:455-547 (converter), :587-604 (datasets), :812-939 (conversion results)
% and :1223-1237 (correspondence ablation), corrected per plan section 1.1 rows 7, 9, 10, 11 and 15.
% Trimmed to 8 pages on 19 Oct: the dataset table T03, the converter table T06, the pipeline figure F02, the v1 stage
% table (tab:fit-stages) and the v2.1 scene table (tab:fit-scenes) were dropped, their numbers kept in the text; the long version is in commit 76aa1a4
% (git show 76aa1a4:docs/thesis/chapters/03_data.tex).
% =============================================================================
\chapter{Data}
\label{ch:data}

This chapter describes the images and labels behind every number in the thesis. Table~\ref{tab:data-splits} lists
every split with its size and its use. I converted 3DPW's SMPL ground truth to Anny so that the model can train and
be tested on real people in Anny space (requirements 1c and 2a). Most of the chapter is about the 3DPW conversion, because the converted labels
decide what the real-image numbers can show. Two rules hold throughout: no test image is trained on, and numbers on
the old 3DPW labels never appear next to numbers on the v23 labels.

\begin{table}[htbp]
\centering\footnotesize
\setlength{\tabcolsep}{4pt}
% Sizes from numbers_data.tex; roles from exp7_final_model.sh:13, :225-226, :741-749 and the sections below.
\begin{tabular}{@{}>{\raggedright\arraybackslash}p{3.0cm}>{\raggedright\arraybackslash}p{2.5cm}>{\raggedright\arraybackslash}p{4.9cm}l@{}}
\toprule
split & size (images / people) & used for & chapters \\
\midrule
AnnyOne training & \DataAnnyOneTrainImages{} images & training of every model; \DataMixAnnyOne\% of each mixed batch & \ref{ch:method}, \ref{ch:results} \\
AnnyOne holdout & \DataHoldoutImages{} images & headline synthetic evaluation (V1); never trained on & \ref{ch:diagnostics}, \ref{ch:results} \\
diagnostic manifest & \DataManifestTrain{} + \DataManifestHoldout{} entries & per-person diagnostics; excluded from every fine-tune & \ref{ch:diagnostics} \\
3DPW-Anny train, v23 & \DataThreeDPWTrainImages{} / \DataThreeDPWTrainPeople & \DataMixThreeDPW\% of each mixed batch & \ref{ch:method}, \ref{ch:results} \\
3DPW-Anny val., v23 & \DataThreeDPWValImages{} / \DataThreeDPWValPeople & \DataValImagesPerEpoch{} images scored after each epoch; monitoring only & \ref{ch:results} \\
3DPW-Anny test, v23 & \DataThreeDPWTestImages{} / \DataThreeDPWTestPeople & reported 3DPW numbers; never trained on & \ref{ch:results} \\
3DPW-Anny old labels & \DataOldTotalImages{} / \DataOldTotalPeople & training of the history models only; retired & \ref{ch:method} \\
Relative Human val. & \DataRHValImages{} / \DataRHValPeople & children and gender on real photos; evaluation only & \ref{ch:diagnostics}, \ref{ch:results} \\
demo images & \DataDemoTotal{} images & by-eye comparison; no metric & \ref{ch:diagnostics}, \ref{ch:results} \\
\bottomrule
\end{tabular}
\caption[Every split and its use]{Every split of the thesis, its size and its use. 3DPW-Anny sizes are images /
labelled people after the drop rules (v23 labels, Section~\ref{sec:data-defects}); numbers on 3DPW are always against
these converted labels, never against the original SMPL meshes. The final model is the last epoch of its run, so the
3DPW validation split does not select it.}
\label{tab:data-splits}
\end{table}

% -----------------------------------------------------------------------------
\section{AnnyOne}
\label{sec:data-annyone}

AnnyOne~\cite{bregier2025anny} is the synthetic dataset released with Anny. It holds \DataAnnyOneImages{} rendered
images of \DataAnnyOneRenderPx{}$\times$\DataAnnyOneRenderPx{} pixels, with \DataAnnyOnePeoplePerImage{} people per
image on average. For every visible person the label gives the exact Anny pose (\DataBones{} bone transforms), the
\DataPhenotypes{} phenotypes and the camera intrinsics. The network sees each image resized to \DataInputPx{}
pixels. The field of view changes from image to image (\DataAnnyOneFovMin{}$^\circ$ to \DataAnnyOneFovMax{}$^\circ$
on a sample of \DataAnnyOneFovSample{} images), so training and evaluation use each image's own intrinsics.

\paragraph{Split.} Indices 0 to \DataAnnyOneLastTrainIndex{} are the training set. The last \DataHoldoutImages{}
images, indices \DataHoldoutFirstIndex{} to \DataHoldoutLastIndex, are the holdout, my own choice and not an
official split; no model trained in this thesis has seen them. NAVER's model was trained on the AnnyOne release, so
for that model the holdout is most likely in-sample. The 3DPW test script also scores the last \DataHoldoutImages{}
AnnyOne images, but in setting V0; its ``AnnyOne test-100'' value is therefore not the holdout number, which is
always given in setting V1.

\paragraph{Few children.} In Anny the child shape appears below an age of about \DataChildAgeLow{} to
\DataChildAgeHigh{} (Section~\ref{sec:bg-bodymodels}). Only \DataAnnyOneAgeBelowThree\% of AnnyOne's people are
younger than 0.3 (\DataAnnyOneAgeBelowTwo\% below 0.2, \DataAnnyOneAgeBelowFour\% below 0.4); the median age is
\DataAnnyOneAgeMedian, and \DataAnnyOneAgeModeShare\% of all people lie between 0.5 and 0.6. In a training batch,
\DataBatchChildLow\% to \DataBatchChildHigh\% of the supervised people are younger than 0.3. The holdout has nobody
younger than \DataHoldoutMinAge{} and tests adults only, which is why I added Relative Human
(Section~\ref{sec:data-rh}).

% -----------------------------------------------------------------------------
\section{3DPW}
\label{sec:data-3dpw}

3DPW~\cite{vonmarcard2018_3dpw} is a real video dataset of \DataThreeDPWSequences{} sequences filmed with a hand-held
phone camera at \DataThreeDPWWidth{}$\times$\DataThreeDPWHeight{} pixels (portrait) and \DataThreeDPWFps{} frames per
second. Inertial sensors on the actors give SMPL ground truth. The standard release has \DataThreeDPWTrainSeq{}
training, \DataThreeDPWValSeq{} validation and \DataThreeDPWTestSeq{} test sequences, with one or two labelled actors
each: \DataThreeDPWPersonSeq{} person-sequences (\DataThreeDPWMen{} men, \DataThreeDPWWomen{} women). Bystanders
have no labels. I keep the splits as 3DPW defines them. One validation sequence, \texttt{outdoors\_crosscountry\_00},
yields no labels (Section~\ref{sec:data-defects}), which leaves \DataThreeDPWValUsableSeq{} usable validation
sequences.

% -----------------------------------------------------------------------------
\section{Converting SMPL to Anny}
\label{sec:data-converter}

The converter (\texttt{smpl\_to\_anny.py}, version 2.3) fits an Anny body to every labelled person, frame by
frame; a label builder moves the fits into 3DPW's camera frame and merges them into one label file per image.

\paragraph{Objective.} Anny has no analytic inverse of its skinning, so I fit by gradient descent. For each person and frame
I search Anny's \DataBones{} rotations $\mathbf{r}$ (6D), the phenotype logits $\boldsymbol{\ell}$ and a translation
$\mathbf{t}$ so that Anny's joints land on the \DataSmplJoints{} SMPL joints $\hat{\mathbf{J}}^{S}$ of the frame:
\begin{equation}
\begin{split}
\min_{\mathbf{r},\boldsymbol{\ell},\mathbf{t}}\
 & \frac{1}{24}\sum_{k=1}^{24}
   \left\| \mathbf{p}_{\pi(k)}\!\left(\mathbf{r},\sigma(\boldsymbol{\ell})\right)
   + \mathbf{t} - \hat{\mathbf{J}}^{S}_{k}\right\|_2^2 \\
 & + \lambda_{r}\sum_{b\neq\text{root}}\left\|\mathbf{R}_b-\mathbf{I}\right\|_F^2
   + \lambda_{\phi}\sum_{j\in\text{free}}\left(\sigma(\ell_j)-\tfrac12\right)^2 ,
\end{split}
\label{eq:fit}
\end{equation}
where $\pi$ is the joint correspondence and $\mathbf{p}_{b}$ the point attached to bone $b$: its head, or a virtual
joint (Section~\ref{sec:data-history}). The first regulariser ($\lambda_r = \DataPoseReg$) pulls every non-root
rotation towards identity, because a twist about a bone's own axis moves no joint but corkscrews the mesh. The second
($\lambda_\phi = \DataShapeReg$) pulls the free phenotypes towards 0.5, so that one remaining joint gap cannot drag a
phenotype to an extreme. Adam (learning rate \DataFitLr) fits the root, then the bone rotations, then the free
phenotypes; shape comes last because it is the weakest signal. Each frame starts from the previous frame's solution
(\DataColdIters{} iterations per stage on a person's first frame, \DataWarmIters{} afterwards). The targets are
3DPW's own per-frame \texttt{jointPositions}; the neutral SMPL model would move them by \DataNeutralShift\,mm on
average, because 3DPW's shape coefficients belong to the gendered models.

\paragraph{Bones and shape.} Only the root and \DataFittedBones{} of the \DataMappedBones{} mapped bones may rotate;
\DataRestBones{} of \DataBones{} bones stay at rest, among them \texttt{neck02}, \texttt{neck03}, \texttt{head} and
every finger bone. Joints cannot determine most phenotypes: weight, muscle, cup size, firmness and the three
ethnicity values move no joint, and adult age barely does. The converter therefore reads gender from 3DPW, fixes age
at \DataFixedAge{} and the seven others at \DataFixedNeutral, and fits only height and proportions. Since version 2.3 it writes men as \DataGenderMale{} and women as
\DataGenderFemale{}, Anny's convention (\texttt{--gender\_convention anny}). After a first pass over all frames, the
median of each free phenotype becomes the person's one body, and every frame's pose is fitted again with it
(\DataRefineIters{} iterations).

\paragraph{Frames, camera and label files.} I fit every third valid frame (\DataLabelFps{} frames per second). Since
version 2.3 one grid holds for every person of a sequence (frame numbers divisible by \DataFrameStride,
\texttt{--frame\_grid shared}), so both actors are fitted on the same images. The builder writes one file per image in the AnnyOne
format and drops a person when the fit error exceeds \DataDropErr\,mm, less than \DataDropVisPct\% of the mesh is
inside the image, or the box is narrower or shorter than \DataDropBox\,px. A dataset class pads the portrait images to a square and shifts the principal
point, so the training and evaluation code run on these labels unchanged.

% -----------------------------------------------------------------------------
\section{Development of the converter}
\label{sec:data-history}

The first converter (v1) chose the Anny bones by name. It gave \DataVOneBarMean\,mm on \texttt{downtown\_bar\_00};
on four test scenes the error was \DataVOneFloorMin{} to \DataVOneFloorMax\,mm on every frame, and all four actors
came out with the same body. A residual independent of the pose and a shape independent of the person mean
that the shape absorbed a fixed mismatch between the skeletons. Placing both models in their rest poses showed two
wrong choices that no name reveals. Anny's rig numbers the spine top-down, SMPL bottom-up, so SMPL's lumbar
\texttt{spine1} had been mapped to a bone \DataSpineMapError\,cm above it; and SMPL's hips had been mapped to
\texttt{pelvis.L/R}, \DataHipMapError\,cm from the hip joint, instead of \texttt{upperleg01.L/R}. Version 2 picks the
spine bones by height along the body axis and the hip bones by distance from the pelvis (Appendix~\ref{app:joints}).
Limbs and head fell to \DataVTwoLimbMin{} to \DataVTwoLimbMax\,mm, but the torso stayed at \DataVTwoTorsoMin{} to
\DataVTwoTorsoMax\,mm (Table~\ref{tab:fit-versions}).

The reason is that the rigs define their torso joints at different anatomical points: SMPL's pelvis joint sits
\DataSmplPelvisAboveHips\,cm above its hips, Anny's root at hip height and \DataAnnyRootFront\,cm in front. For
\DataVirtualJoints{} torso joints I therefore measure once, in the rest pose, the offset $\boldsymbol{\delta}_b$ from
the Anny bone head to the SMPL joint in the bone's own frame, and the fit moves the \emph{virtual joint}
$\mathbf{p}_b = \mathbf{h}_b + \mathbf{R}_b\boldsymbol{\delta}_b$ onto the SMPL joint, with $\mathbf{h}_b$ and
$\mathbf{R}_b$ the posed bone head and frame. The offsets are \DataOffsetPelvis\,cm (pelvis), \DataOffsetHip\,cm
(hips), \DataOffsetSpineMin{} to \DataOffsetSpineMax\,cm (spine) and \DataOffsetCollar\,cm (collars). With them
(v2.1) the torso fell to \DataVTwoOneTorsoMin{} to \DataVTwoOneTorsoMax\,mm and four varied test scenes to
\DataScenesMin{} to \DataScenesMax\,mm (one occluded frame \DataScenesWorstFrame\,mm). Age, however, came out at exactly \DataFreeAge{} for every actor in every
version: joints cannot tell one adult age from another. Version 2.2 fixes age at \DataFixedAge; on all
\DataThreeDPWPersonSeq{} person-sequences it fitted \DataOldFrames{} frames with a median error of
\DataOldFitMedian\,mm. By every measure the converter prints, these fits looked clean. They still held two of the
three defects of the next section.

\begin{table}[htbp]
\centering\small
% Draft main.tex:861-875 (tab:fit_versions), measured; jobs named in smpl_to_anny.py docstring. Cells kept literal.
% Correction: the draft's "fixed gender" for v2 hid that the fixed value was reversed until v2.3.
\begin{tabular}{@{}>{\raggedright\arraybackslash}p{5.2cm}rr>{\raggedright\arraybackslash}p{3.7cm}@{}}
\toprule
converter & mean (mm) & max & worst joints (mm) \\
\midrule
v1: correspondence by name, neutral SMPL targets & \DataVOneFencing & 37 & all joints, common floor \\
v2: correspondence by geometry, 3DPW joints, gender from metadata (reversed until v2.3) & \DataVTwoFencing & 30 & spine1 73, collars 55--59, spine3 58, pelvis 53 \\
v2.1: + virtual torso joints, bone mask, warm start & \DataVTwoOneFencing & 18.8 & shoulder 36, wrists 25--27 \\
\bottomrule
\end{tabular}
\caption[Converter versions on the fencing scene]{Mean joint error over all 24 joints on
\texttt{outdoors\_fencing\_01} (\DataFencingFrames{} frames, including full lunges) for three converter versions.
The geometric correspondence is worth about \DataGainCorrespondence\,mm, the virtual joints a further
\DataGainVirtual\,mm.}
\label{tab:fit-versions}
\end{table}

% -----------------------------------------------------------------------------
\section{Three label defects and the v2.3 re-fit}
\label{sec:data-defects}

The converter grades its own fit, so it cannot see errors in what it is asked to fit or in what happens to the fits
afterwards. I therefore wrote an independent check that reads the finished label files. On sampled images from every
sequence it compares the labels with 3DPW's joints, surface and 2D keypoints, confirms that the loader passes them on
unchanged, and tests gender, a label for every valid actor and the head turn. On the first label set (``old labels'',
v2.2 fits, about \DataOldCheckImages{} images) its grade was FAIL, with three defects (Table~\ref{tab:label-check}).

\paragraph{Defect 1: gender reversed.} The converter wrote 1.0 for men and 0.0 for women, MakeHuman's convention and
the reverse of Anny's, so every 3DPW man was labelled with a female body. Refitting with the flipped value lowered the
joint error for \DataGenderFlipMen{} of \DataGenderFlipMenOf{} men (median \DataGenderFlipMenBefore{} to
\DataGenderFlipMenAfter\,mm); for women it raised the joint error (\DataGenderFlipWomenBefore{} to
\DataGenderFlipWomenAfter\,mm) but brought the surface closer for \DataGenderFlipWomenSurface{} of
\DataGenderFlipWomenOf.

\paragraph{Defect 2: missing actor.} In \DataMissingActor{} of \DataTwoPersonChecked{} sampled two-person images,
one valid actor had no label. Each person's every-third-frame grid started at that person's own first valid frame, so
after an invalid gap the two actors sat on different phases. The detection loss then treats a visible, real person as
background.

\paragraph{Defect 3: missing labels.} The training split had \DataOldTrainImages{} label images where the fits implied
\DataRebuiltTrainImages{} (\DataOldTrainMissingPct\% missing). The image files of \DataMissingImageSeq{} sequences
were incomplete on the cluster (\texttt{courtyard\_captureSelfies\_00}: \DataSelfiesFilesBefore{} of
\DataSelfiesFilesAfter{} files), and the builder skipped frames without an image file without counting them. A
rebuild after re-extracting the images recovered every label (Table~\ref{tab:label-sets}, ``rebuilt''), but still
had defects 1 and 2, so I never trained on it.

\paragraph{The re-fit.} Version 2.3 adds the two switches named above; their defaults reproduce v2.2 exactly. I
re-fitted all \DataThreeDPWPersonSeq{} people with both switches on, and the post-check passed on all \DataRefitSeq{}
sequences (Table~\ref{tab:label-check}, middle block). \DataRefitWorseSeq{} single-person outdoor sequences became
\DataRefitWorseMin{} to \DataRefitWorseMax\,mm worse, likely from a different warm start, which I did not test.

\paragraph{The v23 labels.} I call the labels built from the v2.3 fits the v23 labels: \DataThreeDPWTotalImages{}
images and \DataThreeDPWTotalPeople{} people (Table~\ref{tab:label-sets}). They have fewer images than the rebuilt set
but the same people, because both actors now share each frame. No person was dropped for fit error or box size.
\DataCrossCountryFramesNew{} frames were dropped for visibility, all from \texttt{outdoors\_crosscountry\_00}: with
3DPW's own camera, its ground-truth joints sit at a depth of $\DataCrossCountryDepth$\,m, behind the camera. The 3D
fit is fine; image labels for this sequence cannot exist with this camera, and no label set ever contained it.

% T4 3DPW-Anny label sets (Data 3.5/3.6). Plan REWRITE_PLAN.md section 4, T4.
% Sources: old labels 3575369:31, 3579266:25/:39/:65/:67 (H 3.5); rebuilt v2.2 labels 3581320:25/:39/:65/:67 (H 3.5);
% v23 labels H 0b, H 3.5 "v23 labels built" (builder output, fits of jobs 3590385 + 3602813).
% Layout transposed against the plan (label sets as columns) so that it fits the text width.
\begin{table}[htbp]
\centering\footnotesize
\setlength{\tabcolsep}{4pt}
\begin{tabular}{@{}>{\raggedright\arraybackslash}p{3.6cm}>{\raggedright\arraybackslash}p{2.5cm}>{\raggedright\arraybackslash}p{2.5cm}>{\raggedright\arraybackslash}p{2.5cm}@{}}
\toprule
 & old labels & rebuilt & v23 labels \\
\midrule
converter fits & v2.2 & v2.2 & v2.3 \\
image files of \DataMissingImageSeq{} sequences & missing & re-extracted & re-extracted \\
\addlinespace
train, images / people & \DataOldTrainImages{} / \DataOldTrainPeople & \DataRebuiltTrainImages{} / \DataRebuiltTrainPeople & \DataThreeDPWTrainImages{} / \DataThreeDPWTrainPeople \\
validation & \DataOldValImages{} / \DataOldValPeople & \DataRebuiltValImages{} / \DataRebuiltValPeople & \DataThreeDPWValImages{} / \DataThreeDPWValPeople \\
test & \DataOldTestImages{} / \DataOldTestPeople & \DataRebuiltTestImages{} / \DataRebuiltTestPeople & \DataThreeDPWTestImages{} / \DataThreeDPWTestPeople \\
total & \DataOldTotalImages{} / \DataOldTotalPeople & \DataRebuiltTotalImages{} / \DataRebuiltTotalPeople & \DataThreeDPWTotalImages{} / \DataThreeDPWTotalPeople \\
\addlinespace
gender reversed & yes & yes & no \\
actor missing in two-person frames & yes & yes & no \\
labels missing (no image file) & yes & no & no \\
used for training & history models only & never & Final, Control, Student \\
\bottomrule
\end{tabular}
\caption{The three 3DPW-Anny label sets. All are built at every third frame with the same drop rules. The v23
labels have fewer images than the rebuilt set but the same people, because both actors of a two-person
sequence now share each frame. \texttt{outdoors\_crosscountry\_00} yields no labels in any set
(Section~\ref{sec:data-defects}). Numbers measured on the old labels are retired and never appear next to
numbers on the v23 labels.}
\label{tab:label-sets}
\end{table}


\paragraph{The check on the v23 labels.} The same check on the v23 labels (\DataCheckPeople{} people in
\DataCheckImages{} images, \DataCheckSeq{} sequences) moved the grade from FAIL to WARN
(Table~\ref{tab:label-check}). The median body-joint error fell from \DataCheckJointsOld{} to
\DataCheckJointsNew\,mm, flipping the gender now makes the fit worse, and no sampled two-person image misses an actor.
The head turn relative to the chest is now kept; in the old test labels it was mostly lost
(\DataCheckHeadTestLabelOld$^\circ$ against \DataCheckHeadTestSmplOld$^\circ$ in SMPL). Two rows stay at WARN, the
limb-length ratios and the people with a joint more than \DataCheckFarMm\,mm off, mostly at the hip, likely because
Anny's and SMPL's limb proportions and hip joints differ (Figure~\ref{fig:label-check}).

% T5 Label quality before and after the re-fit (Data 3.5). Plan REWRITE_PLAN.md section 4, T5; H 0b table; H 3.5.
% Columns: old labels = check job 3575369 (15 Oct, v2.2 fits); v23 labels = check job 3611284 (18 Oct, v2.3 fits).
% Re-fit post-check rows: logs/check_refit_all_v23.log (H 3.5 "Re-fit v2.3 result").
\begin{table}[htbp]
\centering\footnotesize
\setlength{\tabcolsep}{4pt}
\begin{tabular}{@{}>{\raggedright\arraybackslash}p{5.6cm}>{\raggedright\arraybackslash}p{2.9cm}>{\raggedright\arraybackslash}p{3.1cm}@{}}
\toprule
check & old labels & v23 labels \\
\midrule
overall grade of the check & FAIL & WARN \\
body joints vs 3DPW, median / p95 (mm) & \DataCheckJointsOld{} / \DataCheckJointsOldPn & \DataCheckJointsNew{} / \DataCheckJointsNewPn \\
surface, Anny$\to$SMPL / SMPL$\to$Anny (mm) & \DataCheckSurfAOld{} / \DataCheckSurfBOld & \DataCheckSurfANew{} / \DataCheckSurfBNew \\
2D error over raw 3DPW (px) & $\DataCheckTwoDOld$ & $\DataCheckTwoDNew$ \\
gender labels & reversed & men 0, women 1 \\
\quad refit with flipped gender & better for \DataGenderFlipMen{} of \DataGenderFlipMenOf{} men & worse: men \DataCheckFlipMenLabel{} $\to$ \DataCheckFlipMenFlipped\,mm \\
two-person images with an unlabelled actor & \DataMissingActor{} of \DataTwoPersonChecked & 0 of \DataCheckGapNewOf \\
fitted frames; per-frame error, median (mm) & \DataOldFrames; \DataOldFitMedian & \DataCheckFramesNew; \DataCheckFitMedianNew \\
head turn vs chest, labels / SMPL, median (deg) & \DataCheckHeadLabelOld{} / \DataCheckHeadSmplOld & \DataCheckHeadLabelNew{} / \DataCheckHeadSmplNew \\
\midrule
\multicolumn{3}{@{}l}{\emph{post-check of the re-fit, all fitted frames, \DataThreeDPWPersonSeq{} people}} \\
mean joint error, men (\DataThreeDPWMen) / women (\DataThreeDPWWomen) (mm) & \DataRefitMenBefore{} / \DataRefitWomenBefore & \DataRefitMenAfter{} / \DataRefitWomenAfter \\
two-person frames with one valid actor unfitted & \DataRefitGapBefore{} of \DataRefitGapBeforeOf & 0 of \DataRefitGapAfterOf \\
\midrule
\multicolumn{3}{@{}l}{\emph{still WARN on the v23 labels}} \\
limb-length ratio Anny / SMPL, forearm / thigh & & \DataCheckForearm{} / \DataCheckThigh \\
people with a joint more than \DataCheckFarMm\,mm off & & \DataCheckFarPct\% \\
\bottomrule
\end{tabular}
\caption[Conversion check before and after the re-fit]{Independent conversion check of the 3DPW-Anny labels before (old labels, v2.2 fits) and after the
re-fit (v23 labels, v2.3 fits). The check samples images from every sequence and compares the labels with
3DPW's own joints, surface and 2D keypoints. The middle block is the re-fit's own post-check over all fitted
frames; old and new fits fitted different frames, so it is not a paired comparison.}
\label{tab:label-check}
\end{table}


\begin{figure}[tbp]
  \centering
  \IfFileExists{figures/img/F14_label_check.jpg}{%
    \includegraphics[width=\linewidth]{figures/img/F14_label_check.jpg}%
  }{%
    \fbox{\parbox{0.9\linewidth}{\centering\vspace{2em}Figure pending: \texttt{F14\_label\_check.jpg}, three
    panels from the conversion check of the v23 labels, one of them a two-person frame.\vspace{2em}}}%
  }
  \caption[Conversion check on the v23 labels]{Conversion check on the v23 labels: the 3DPW photo, 3DPW's SMPL ground truth with its 2D keypoints, and
  our Anny label drawn with 3DPW's camera. Real 3DPW frames. By design the labels have rest-pose hands and a head
  that turns only through \texttt{neck01} (Section~\ref{sec:data-limits}).}
  \label{fig:label-check}
\end{figure}

% -----------------------------------------------------------------------------
\section{What the 3DPW labels can and cannot test}
\label{sec:data-limits}

The v23 labels are converted labels, not ground truth. Four limits follow from the converter's design.

\paragraph{Eight constant phenotypes.} The labels test three phenotypes: gender (from metadata), height and
proportions (fitted; medians \DataLabelMedianHeight{} and \DataLabelMedianProportions). The other
\DataConstantCount{} are the same for every person: age \DataFixedAge{} and \DataNeutralCount{} values at
\DataFixedNeutral. A model trained on 3DPW is pulled towards these constants for every real person, by the shape loss
and the mesh terms; 3DPW people are about \DataSupervisedShareThreeDPW\% of the supervised people in a 70/30 batch
(my estimate). The 3DPW error is also measured against these neutral-body meshes, so it likely rewards a neutral
prediction; I did not test this. Section~\ref{sec:method-real} describes the shape mask that removes the shape loss on
the constants (Final-masked). At age \DataFixedAge{} every label is a full adult
(Figure~\ref{fig:bg-age-sweep}); 3DPW has no children, so it can neither teach nor test them.

\paragraph{Head and hands at rest.} \texttt{neck02}, \texttt{neck03} and \texttt{head} stay at identity in every
label, so the head turns only through \texttt{neck01}, and the labels cannot test the half-turned neck bones of
Chapter~\ref{ch:diagnostics}. 3DPW's SMPL ground truth has no finger joints, so every label has rest-pose hands; the
training masks their rotation terms for 3DPW people (Section~\ref{sec:method-real}).

\paragraph{A floor of about \DataFloor\,mm.} The labels sit \DataCheckJointsNew\,mm (median) and
\DataCheckJointsNewPn\,mm (p95) from 3DPW's own joints. Every Anny-space error on 3DPW contains this residual.

\paragraph{Not comparable with published numbers.} I never evaluated against the original SMPL meshes, and our MPJPE
averages all \DataBones{} Anny joints (Section~\ref{sec:method-eval}), so I write ``lowest error on our 3DPW-Anny
labels''. Bystanders are unlabelled, so precision on 3DPW is a lower bound.

% -----------------------------------------------------------------------------
\section{Relative Human}
\label{sec:data-rh}

Neither the holdout nor 3DPW contains children. Relative Human, introduced with BEV~\cite{sun2022bev}, fills this gap.
It holds \DataRHImages{} real photos, many crowded and with children. Each person has a box, an age group (adult,
teen, kid, baby) and a gender coded 0 for male, as in Anny; there are no 3D labels. The dataset ships without a
licence file, so I use only the validation split and only for evaluation: \DataRHValImages{} images with
\DataRHValPeople{} labelled people (\DataRHValAdults{} adults, \DataRHValTeens{} teens, \DataRHValKids{} kids,
\DataRHValBabies{} babies, \DataRHValUnknownAge{} of unknown age).

The model runs once on each whole photo with a field of view of \DataRHFov$^\circ$, and detections are paired with
labelled boxes as in Section~\ref{sec:method-eval}. A predicted age below \DataRHChildThr{} counts as child-shaped, a
predicted gender below \DataRHGenderThr{} as male. I report recall and mean predicted age per age group, the
child-against-adult AUC of the predicted age (kids and babies against adults) and the gender accuracy. Every number
covers detected people only, so recall always stands next to it.

% -----------------------------------------------------------------------------
\section{The diagnostic manifest and the demo images}
\label{sec:data-manifest}

The checks of Chapter~\ref{ch:diagnostics} need the same people every time. A seeded script picks them from the
AnnyOne labels alone and stores \DataManifestTotal{} fixed (image, person) entries: \DataManifestTrain{} from the
training range (adults, the youngest children, clearly visible hands, the hardest head turns) and
\DataManifestHoldout{} from the holdout. Base saw the training entries; every later fine-tune gives their images a
sampling weight of zero, so they remain the same in-sample reference for all models. The manifest children therefore
show whether a model can fit a child it has seen, not whether it generalises. For by-eye comparisons I use
\DataDemoTotal{} real images: \DataDemoPhotos{} photos and \DataDemoFrames{} 3DPW test frames, one middle frame from
each of \DataDemoFrames{} test sequences. No model trained on them and no metric is computed on them.
